\begin{landscape}

\chapter{Orçamento}

\begin{table}[htpb]
\begin{center}
\caption{Orçamento geral.}
\begin{tabular}{|c|c|c|}
\hline
\textbf{Geral}          & \textbf{Previsto (R\$)} & \textbf{Realizado (R\$)} \\ \hline
Matéria-prima           & 1,00                    & 1,00                     \\ \hline
Componentes eletrônicos & 2,00                    & 3,00                     \\ \hline
Componentes elétricos   & 2,00                    & 2,50                     \\ \hline
Ferramentas             & 3,00                    & 2,75                     \\ \hline
Montagem                & 4,00                    & 2,00                     \\ \hline
\end{tabular}
\end{center}
\end{table}

\begin{table}[htpb]
\begin{center}
\caption{Orçamento por aquisição.}
\begin{tabular}{|c|c|c|c|c|}
\hline
\textbf{Item ou serviço} & \textbf{Tipo} & \textbf{Previsto (R\$)} & \textbf{Realizado (R\$)} & \textbf{Responsável} \\ \hline
Madeira      & Matéria-prima & 0,50   & 0,60 & Fulano                    \\ \hline
Ferro        & Matéria-prima & 0,25   & 0,25 & Fulano                    \\ \hline
Plástico ABS & Matéria-prima & 0,25   & 0,15 & Fulano                    \\ \hline
Arduino      & Componente eletrônico & 2,00  & 3,00 & Beltrano                     \\ \hline
Fonte DC     & Componente elétrico & 2,00   & 2,50 & Beltrano                     \\ \hline
Martelo & Ferramentas & 3,00   & 2,75 & Fulano                     \\ \hline
Impressão 3D & Montagem & 4,00   & 2,00 & Ciclano                   \\ \hline
\end{tabular}
\end{center}
\end{table}

\begin{table}[htpb]
\begin{center}
\caption{Justificativas para itens ou serviços trocados após a aquisição.}
\begin{tabular}{|c|c|} \hline
\textbf{Item ou serviço} & \textbf{Justificativa}  \\ \hline 
    Martelo &
    O material solicitado não foi entregue no prazo pelo fornecedor. \\ \hline
    Ferro &
    Erro no dimensionamento do primeiro material. \\ \hline
\end{tabular}
\end{center}
\end{table}

\section{Orçamento - Estrutura}

\noindent O orçamento da área de Estrutura contempla tanto os materiais e insumos necessários para a construção do labirinto quanto os custos associados à fabricação do chassi do robô. Para o labirinto, foram considerados os gastos com MDF para a base e para as paredes, tintas para acabamento, além de parafusos, buchas e pinos cavilha para fixação e montagem. Para o chassi, foi estimado um consumo aproximado de 50 g de filamento, com base em um valor de R\$ 150,00 por quilograma, resultando em R\$ 7,50 de material. Além disso, foi previsto um valor de R\$ 40,00 para o serviço de impressão 3D, considerando uso do equipamento, preparação do arquivo, tempo de impressão, energia, desgaste da máquina e mão de obra, bem como R\$ 10,00 destinados aos parafusos de fixação. Dessa forma, o orçamento total da área de Estrutura foi estimado em R\$ 167,50.

\begin{table}[htpb]
\begin{center}
\caption{Orçamento geral atualizado da Estrutura.}
\begin{tabular}{|c|c|c|}
\hline
\textbf{Geral}          & \textbf{Previsto (R\$)} & \textbf{Realizado (R\$)} \\ \hline
Matéria-prima           & 107,50                  & ---                      \\ \hline
Fixadores               & 18,00                   & ---                      \\ \hline
Acabamento              & 52,00                   & ---                      \\ \hline
Montagem                & 40,00                   & ---                      \\ \hline
\textbf{Total}          & \textbf{167,50}         & \textbf{---}             \\ \hline
\end{tabular}
\end{center}
\end{table}

\noindent A divisão adotada neste orçamento considera como matéria-prima o MDF do labirinto e o PLA do chassi. Os custos com fixadores incluem parafusos, buchas, pinos cavilha e parafusos do chassi. O acabamento corresponde às tintas utilizadas no labirinto, enquanto a montagem contempla o serviço de impressão 3D.

\begin{table}[htpb]
\begin{center}
\caption{Orçamento por aquisição atualizado da Estrutura.}
\begin{tabular}{|p{7cm}|p{3.3cm}|c|c|c|}
\hline
\textbf{Item ou serviço} & \textbf{Tipo} & \textbf{Previsto (R\$)} & \textbf{Realizado (R\$)} & \textbf{Responsável} \\ \hline
MDF para base do labirinto          & Matéria-prima & 26,00  & --- & Estrutura \\ \hline
MDF para paredes do labirinto       & Matéria-prima & 12,00  & --- & Estrutura \\ \hline
Tintas para acabamento do labirinto & Acabamento    & 52,00  & --- & Estrutura \\ \hline
Parafusos e buchas do labirinto     & Fixadores     & 12,00  & --- & Estrutura \\ \hline
Pinos cavilha do labirinto          & Fixadores     & 8,00   & --- & Estrutura \\ \hline
Filamento PLA para chassi           & Matéria-prima & 7,50   & --- & Estrutura \\ \hline
Serviço de impressão 3D do chassi   & Montagem      & 40,00  & --- & Estrutura \\ \hline
Parafusos do chassi                 & Fixadores     & 10,00  & --- & Estrutura \\ \hline
\textbf{Total}                      & ---           & \textbf{167,50} & \textbf{---} & --- \\ \hline
\end{tabular}
\end{center}
\end{table}

\begin{table}[htpb]
\begin{center}
\caption{Justificativas da Estrutura para itens ou serviços trocados após a aquisição.}
\begin{tabular}{|p{5cm}|p{16cm}|} \hline
\textbf{Item ou serviço} & \textbf{Justificativa}  \\ \hline
    --- &
    Até o momento, não houve substituição de itens ou serviços após a aquisição, pois a seção apresenta apenas os valores previstos. \\ \hline
\end{tabular}
\end{center}
\end{table}

\end{landscape}
